% GENERATED by md2tex.py from ../draft_v5_en/back/B_appendix_requirements.md -- do not edit by hand.
% Change the Markdown and run `python md2tex.py`; edits here are overwritten.

\chapter[Safety requirements specification and rationale]{Safety requirements specification and \emph{rationale}}
\label{app:B}

This is the full version of Table~\ref{tab:4.2}, with thresholds and traceability. It is followed by the \emph{rationale} for each requirement: where each threshold comes from, what makes it falsifiable and what state its calibration is in.

% layout: longtable, fixed, \footnotesize, 164 of 165 mm
\begingroup
\footnotesize\setlength{\tabcolsep}{4pt}
\setlength{\LTleft}{0pt}\setlength{\LTright}{\fill}\setlength{\LTcapwidth}{\textwidth}
\begin{longtable}{@{}>{\RaggedRight\arraybackslash}p{9.7mm}>{\RaggedRight\arraybackslash}p{48.1mm}>{\RaggedRight\arraybackslash}p{22.2mm}>{\RaggedRight\arraybackslash}p{13.3mm}>{\RaggedRight\arraybackslash}p{13.3mm}>{\RaggedRight\arraybackslash}p{19.3mm}>{\RaggedRight\arraybackslash}p{20.9mm}@{}}
\caption{Safety requirements with thresholds and traceability (full version of Table~\ref{tab:4.2}).}\label{tab:B.1} \\
  \toprule
  \tabhead{ID} & \tabhead{Short statement} & \tabhead{Main parameters} & \tabhead{H covered} & \tabhead{Cage rule} & \tabhead{Criticality} & \tabhead{Verification} \\
  \midrule
  \endfirsthead
  \multicolumn{7}{@{}l}{\emph{(continued)}} \\
  \toprule
  \tabhead{ID} & \tabhead{Short statement} & \tabhead{Main parameters} & \tabhead{H covered} & \tabhead{Cage rule} & \tabhead{Criticality} & \tabhead{Verification} \\
  \midrule
  \endhead
  \midrule
  \multicolumn{7}{r@{}}{\emph{(continued on next page)}} \\
  \endfoot
  \bottomrule
  \endlastfoot
  \sr{001} & The absolute lateral offset with respect to the centre of the road shall stay below \texttt{d\_max} during operation inside the applicable ODD. & \texttt{d\_max = 0.16 m} & \hz{01} & \cagerule{01} & \mbox{SR-CL-A} & \scn{NOM-01}, \scn{NOM-02}, \scn{EDGE-02} \\
  \addlinespace[2pt]
  \sr{002} & The absolute orientation error with respect to the lane direction shall stay below \texttt{θ\_max}. & \texttt{θ\_max = 0.44 rad (25°)} & \hz{02} & \cagerule{02} & \mbox{SR-CL-A} & \scn{EDGE-01}, \scn{EDGE-04} \\
  \addlinespace[2pt]
  \sr{003} & The projected time to lane crossing (TTLC) shall stay above \texttt{t\_min}; the 5th percentile of TTLC over nominal runs shall be $\ge$ 0.5 s. & \texttt{t\_min = 1.0 s}; p5 floor = 0.5 s & \hz{01}, \hz{02} (partial) & \cagerule{03} & \mbox{SR-CL-A} & \scn{NOM-02}, \scn{EDGE-01} \\
  \addlinespace[2pt]
  \sr{004} & The longitudinal speed shall not exceed \texttt{v\_max(κ)}, a ceiling dependent on the local curvature $\kappa$. & \texttt{v\_max\_straight = 0.5 m/s}; \texttt{v\_max\_curve = 0.25 m/s} & \hz{03} & \cagerule{04} & \mbox{SR-CL-A} & \scn{NOM-02}, \scn{EDGE-03} \\
  \addlinespace[2pt]
  \sr{005} & Under a compound trigger on heading and offset lasting \texttt{Δt\_max}, the system shall transition to emergency mode with minimum deceleration and frozen steering. & \texttt{θ\_warn = 20°}; \texttt{d\_warn = 0.12 m}; \texttt{Δt\_max = 0.2 s}; \texttt{a\_min = 0.3 m/s²} (provisional, subject to \mbox{M-3}) & \hz{04}, \hz{07} (partial) & \cagerule{05} & \mbox{SR-CL-A} & \scn{EDGE-04} \\
  \addlinespace[2pt]
  \sr{006} & The command variation between two consecutive cycles shall stay below \texttt{δ\_max} for steering and throttle. & \texttt{δ\_max\_steer = 0.15}; \texttt{δ\_max\_thr = 0.10} (per cycle) & \hz{05} & \cagerule{06} & \mbox{SR-CL-B} & All scenarios (rate limiter active) \\
  \addlinespace[2pt]
  \sr{007} & The cage shall activate emergency mode if the observation is older than \texttt{staleness\_max} or if any field is outside the plausible range. & \texttt{staleness\_max = 200 ms}; \texttt{N\_missing\_max = 5 cycles} & \hz{06} & part of \cagerule{05} & \mbox{SR-CL-A} & \scn{PERT-02} \\
  \addlinespace[2pt]
  \sr{008} & Under an external stop signal or a controlled end of episode, the system shall decelerate to 0 m/s within \texttt{t\_stop\_max} without exceeding the lateral \texttt{d\_max}. & \texttt{t\_stop\_max = 1.7 s}; \texttt{d\_max = 0.16 m} & \hz{07} & part of \cagerule{05} + vehicle-control node & \mbox{SR-CL-A} & \scn{NOM-03}, \scn{EDGE-04} \\
  \addlinespace[2pt]
  \sr{009} & In every eligible window of \texttt{t\_window}, the vehicle shall accumulate at least \texttt{Δs\_min} of nominal longitudinal progress; the metric \met{S2} under monitoring mode shall not exhibit a sustained rise with respect to the baseline. & \texttt{Δs\_min = 0.10 m}; \texttt{t\_window = 2.0 s}; \texttt{Δt\_settle = 1.0 s} & \hz{08} & training (\dec{25}) & \mbox{SR-CL-B} & \scn{NOM-01}..03, \scn{PERT-03} \\
  \addlinespace[2pt]
  \sr{010} & When two or more cage rules activate in the same cycle, the final command shall satisfy the safe envelope of every activated rule and the inter-cycle pattern shall not exhibit sustained oscillation above \texttt{f\_osc\_max}. & \texttt{f\_osc\_max = 5 Hz}; joint-envelope assertion & \hz{09} & arbiter (\dec{25}) & \mbox{SR-CL-B} & \scn{EDGE-04}, \scn{EDGE-05} \\
  \addlinespace[2pt]
  \sr{011} & The standard deviation of the heading error over eligible windows of \texttt{t\_psd} shall stay below \texttt{σ\_θ\_max}, covering the in-band branch of \hz{02} that \sr{002} does not bound. & \texttt{σ\_θ\_max = 5°}; \texttt{t\_psd = 1.0 s} & \hz{02} (oscillatory branch) & \cagerule{06} + training & \mbox{SR-CL-B} & \scn{EDGE-01}, \scn{EDGE-04} \\
  \addlinespace[2pt]
  \sr{012} & (Track \enquote*{E}) Lane keeping shall stay inside the \sr{001}/\allowbreak{}\sr{002} envelope under degraded visual input (glare, exposure, motion blur, contrast, shadow). & \texttt{d\_max = 0.16 m}; \texttt{θ\_max = 25°} (reused); provisional visual envelope (docs/09) & \hz{10} & \cagerule{01}, \cagerule{02}, \cagerule{03} + training & \mbox{SR-CL-A} & \scn{PERT-04}, \scn{PERT-05}, \scn{PERT-06}, \scn{PERT-09}, \scn{PERT-10} \\
  \addlinespace[2pt]
  \sr{013} & (Track \enquote*{E}) On loss of valid perception (occlusion / absence of features / stale or dropped frames), the system shall enter an open-loop controlled stop through \cagerule{05} without exceeding \texttt{d\_max}. & \texttt{perc\_\allowbreak{}staleness\_\allowbreak{}max = 200 ms} (provisional); \texttt{d\_max = 0.16 m} & \hz{11} & part of \cagerule{05} (health of the CV estimator of the cage) & \mbox{SR-CL-A} & \scn{PERT-07} \\
  \addlinespace[2pt]
  \sr{014} & (Track \enquote*{E}) The cage shall not impose its rules on a lane estimate that fails the plausibility / temporal consistency check; instead it shall enter a controlled stop (\cagerule{05}). & \texttt{plaus\_tol}, \texttt{Δt\_plaus} (provisional, against the ground-truth oracle) & \hz{12} & part of \cagerule{05} (plausibility check $\rightarrow$ stop) & \mbox{SR-CL-A} & \scn{PERT-08}, \scn{PERT-04}..06, \scn{PERT-09}, \scn{PERT-10} \\
\end{longtable}
\endgroup

\section{Rationale per requirement}
\label{sec:B.1}

The full rationale for each SR is in \texttt{docs/\allowbreak{}03\_\allowbreak{}safety\_\allowbreak{}requirements.\allowbreak{}md}. It includes the physical justification of the threshold based on the ODD parameters, the falsifiable form of the requirement, the experiment that verifies it and the reference to the hazard it comes from. To show the level of detail required, this subsection gives the full rationale of \sr{001} as an example. The rest is summarised at the end and cited by reference.

The rationale of \sr{001} is as follows. The parameter \texttt{d\_max = 0.16 m} is not chosen based on the policy's performance, but on the geometric limits of the ODD. The road has a total width of \texttt{ODD-1.ROAD\_WIDTH = 0.50 m}, so the physical edge of the drivable corridor is \texttt{0.25 m} from the track axis. A total safety margin of \texttt{Δ = 0.09 m} covers three independent contributions: the lateral noise of the state estimator (\texttt{≈ 0.01 m}), the maximum drift expected from the nominal control latency (\texttt{v\_max · LATENCY\_NOMINAL = 0.5 m/s · 50 ms = 0.025 m}) and half the physical width of the 1:14 CobraFlex (\texttt{≈ 0.05 m}). The falsifiable threshold is then \texttt{d\_max = ROAD\_WIDTH/2 − Δ = 0.25 − 0.09 = 0.16 m}, read as the absolute lateral offset of the vehicle's geometric centre from the axis of the drivable corridor. This sign and unit convention is fixed in the statement of \sr{001} in \texttt{docs/\allowbreak{}03\_\allowbreak{}safety\_\allowbreak{}requirements.\allowbreak{}md}, and the SRs that depend on it cite it (\sr{005}, with \texttt{d\_warning = 0.12 m < d\_max} as an early warning threshold, and \sr{008}, which requires staying below \texttt{d\_max} during the controlled stop).

The rationale for the other SRs is summarised below, with a pointer to the separate document for the details.

\textbf{\sr{002}} (\texttt{θ\_max = 25°}) is justified with a recoverability calculation based on the bicycle model, with a wheelbase of 0.15 m, steering saturated at 0.5 rad, a nominal speed of 0.3 m/s and a cage response time of 0.05 s. The value falls inside the recoverable region with a margin of roughly a factor of two.

\textbf{\sr{003}} (\texttt{t\_min = 1.0 s}) is split into 0.3 s of margin for the cage (justifiable with kinematics) and 0.7 s of margin for the policy (marked provisional, to be reviewed after the F3 training prototype).

\textbf{\sr{004}} defines a speed ceiling that depends on curvature, with \texttt{ODD-1.V\_MAX = 0.5 m/s} on straights and 0.25 m/s in curves. The interpolation coefficient \texttt{k\_κ = 0.3} is chosen so that the ceiling drops exactly at the maximum curvature expected on the \texttt{odd3\_curvy\_loop} map (pending closure with \mbox{TBD-Q9}). \textbf{\sr{005}} introduces a compound trigger with a persistence of \texttt{Δt\_max = 0.2 s} (four control cycles, needed to tell a real compound state apart from a passing glitch). \texttt{a\_min = 0.3 m/s²} is provisional until the \mbox{M-3} measurement on the platform.

\textbf{\sr{006}} sets a rate limiter as a cautious defence against abrupt actuation. The values \texttt{δ\_max\_steer = 0.15} and \texttt{δ\_max\_thr = 0.10} are defaults that still need to be checked against the mechanical limits of the actuator (measurement \mbox{M-5}) and against the 95th percentile of the trained policy's natural command change (after the F3 prototype).

\textbf{\sr{007}} keeps a staleness limit of \texttt{staleness\_max = 200 ms} (four control cycles) together with a counter of missing messages, \texttt{N\_missing\_max = 5}. Both values are cautious compared with the nominal properties of the ROS2 bus. The plausible ranges for each state field are kept wider than the operating limits of each variable on purpose, so that an out-of-range value is a clear sign of sensor failure.

\textbf{\sr{008}} sets \texttt{t\_stop\_max = 1.7 s}, consistent with \texttt{v\_max\_straight / a\_min ≈ 1.67 s} from \sr{005} plus a margin for granularity and latency. The fix of the earlier mismatch with \sr{005} (1.5 s in the F0 baseline) is recorded in \texttt{docs/\allowbreak{}CHANGELOG.\allowbreak{}md}.

\textbf{\sr{009}} sets a minimum on longitudinal progress. \texttt{Δs\_min = 0.10 m} is the product of the minimum useful operating speed (\texttt{v\_min ≈ 0.05 m/s}) and a sliding window of \texttt{t\_window = 2.0 s}. The exception of \texttt{Δt\_settle = 1.0 s} after leaving emergency mode or a controlled stop keeps \sr{009} from clashing with \sr{005} or \sr{008} while the vehicle starts moving again (see §\sr{009} of the SRS for the explicit priority order). It is implemented in training (\dec{25}): the cage does not force progress, it only detects the stall through \met{P6} and signals the test harness.

\textbf{\sr{010}} makes sure the cage combines its rules consistently. The \emph{joint-envelope assertion} at the end of each cycle checks that the command sent out meets the preconditions of every rule active in that cycle, and a monitor across cycles keeps oscillation between contradictory corrections below \texttt{f\_osc\_max = 5 Hz}. It is not an extra numbered rule but a structural property of the cage pipeline (\texttt{arbiter}, \dec{25}). If the assertion fails, the system drops into \cagerule{05} (emergency mode) as a last line of defence.

\textbf{\sr{011}} covers the oscillating version of \hz{02}, which the magnitude threshold of \sr{002} does not limit. \texttt{σ\_θ\_max = 5°} allows oscillations up to about 7° of amplitude, with plenty of margin against \texttt{θ\_max = 25°}, but it is strict enough to catch oscillation that stays within bounds. The window of \texttt{t\_psd = 1.0 s} covers at least one period of the relevant oscillation ($\approx$ 1 Hz) without mixing it with slower drifts. The implementation is mixed (\texttt{C-06 + training}): \cagerule{06} smooths the high-frequency content of the policy's commands, and a heading variance term in the reward encourages the policy not to oscillate in steady state.
